% sections/01-introduction.tex
\section{Introduction}
\label{sec:intro}

Between 2024 and 2026, vulnerability discovery, exploitation, and patching acquired a new
subject: software systems in which large language models select and execute security-relevant
actions. Google reported an agent finding an exploitable memory-safety zero-day in SQLite that
150 CPU-hours of AFL could not rediscover~\cite{bigsleep2024naptime}; a commercial agent
briefly held the top US position on a major bug-bounty platform~\cite{xbowtop12025}; a
government challenge scored seven fully autonomous systems on finding, proving, and patching
vulnerabilities in open-source software~\cite{darparesults2025,r7sokaixcc2026}; and Anthropic
reported an intrusion campaign whose execution was 80--90\% autonomous~\cite{anthropicgtg2025}.
These events are connected neither by a shared benchmark nor by a shared metric, which is
precisely why existing surveys cannot answer the questions they raise.

\subsection{The Empirical Problem}
Three difficulties block straightforward assessment. First, \emph{unit confusion}: results are
reported about models (benchmark pass rates), pipelines (fuzz-target generation), agents
(tool-using loops), and organizations (leaderboard ranks) as if commensurable. Second,
\emph{loop fragmentation}: discovery, validation, and repair are evaluated separately even when
the deployed systems integrate them. Third, \emph{validity opacity}: headline numbers travel
without the information conditions, label provenance, or adjudication rules that produced them
--- GPT-4-based agents exploited 87\% of one-day CVEs when given their descriptions and 7\%
when not~\cite{fang2024agents}, yet the qualified number is the one usually quoted.

\subsection{Our Approach}
We build a PRISMA-grounded, multi-stream evidence base (academic literature; competition and
industrial documents; registry records) and classify every record on two orthogonal axes:
a \emph{decision-authority} gradient A0--A3, assigned strictly by documented behavior ---
``who chooses the next action?'' --- with autonomy claims graded by their weakest missing
condition; and the offense/defense loop positions B1--B4, recorded as sets for systems spanning
phases. Fourteen in-window records bear autonomy classifications; the remaining 22 provide
context, validity evidence, loop-transition documentation, or related-work positioning, and we
keep the distinction explicit throughout.

\subsection{Contributions}
\textbf{C1.} An autonomy$\times$loop taxonomy applied record-by-record with published
behavioral evidence (Secs.~\ref{sec:framework}, \ref{sec:discovery}; Fig.~\ref{fig:abmap}).
\textbf{C2.} An evidence-graded analysis of offense$\leftrightarrow$defense coupling, including
a documented cross-program causal edge from the AIxCC competition to Project Zero's research
program (Sec.~\ref{sec:coevolution}). \textbf{C3.} A systematization of evaluation-validity
failures as measurable phenomena, with a five-part contamination-control protocol
(Sec.~\ref{sec:validity}). \textbf{C4.} A falsifier-specified experimental roadmap
(Sec.~\ref{sec:roadmap}).

\subsection{Insights}
\textbf{I1.} Evaluation, deployment, and discourse increasingly measure progress in terms of
agentic autonomy rather than model capability alone; whether this reflects a corresponding
shift in underlying capability remains unresolved. \textbf{I2.} Multiple documented
offense$\leftrightarrow$defense coupling edges exist --- including one explicit cross-program
inspiration --- but the evidence does not establish field-wide causal co-evolution.
\textbf{I3.} Evaluation validity is an increasingly binding constraint on what this field can
claim to know.

\subsection{Scope}
We cover January 2024--June 2026, with pre-2024 work as background only. We classify
\emph{documented behavior}, not marketing; we distinguish vendor-reported from independently
verified results everywhere; and we found it necessary to state plainly that no included
evidence met our operational definition of an autonomous production hunter (A3).
